% =====================================================================
\section{The tools and their free-field acid tests}

\figfredholm
\label{ch:tools}
% =====================================================================

\subsection{The acid test}

The project charter imposes a \textbf{free-field acid test} on every purported new tool: it must
fail on the free massless field, where $f_1^2=0$ and $G''=0$. A tool that ``proves'' a gap for
the free field is vacuous and must be discarded. The test is applied in three adversary classes,
in increasing order of severity:
\begin{enumerate}
\item \textbf{Free massless field}: $\nu$ supported on $[0,\infty)$ with no atom at zero,
  $f_1^2=0$, $G''=0$. The tool must return $0$ or a non-positive value.
\item \textbf{Composite $\mathord{:}\varphi^2\mathord{:}$}: a gapped theory whose $0^{++}$
  operator is composite; $f_1^2=0$ for the continuum $\mathord{:}\varphi^2\mathord{:}$ channel at
  the bottom (the state is a two-particle threshold), and $G''=1/\chi\ne m_1^2$. The tool must
  also return $0$ here; a tool that confuses ``gapped'' with ``atom'' fails this class.
\item \textbf{Free gapped field}: a massive free scalar, which has a gap but no $0^{++}$ atom in
  the chosen channel. The tool must not mistake the gap for an atom.
\end{enumerate}
A tool passes if it is correct on all three classes and discriminatingly positive on
Yang--Mills. ``Discriminatingly'' means: not merely positive on YM, but negative or zero exactly
where it should be, so that the positivity carries information.

\subsection{The six first-wave tools}
\label{sec:sixtools}

\subsubsection{Multiplicative residue transport}

\textbf{Formula.}
$f_1^2(\mu)=Z_O(\mu,\mu_0)^2\,f_1^2(\mu_0)$, with
$Z_O=(\beta/2g)(\mu_0)/(\beta/2g)(\mu)>0$ and anomalous dimension
$\gamma_O=-d\log(\beta/2g)/d\log\mu=2b_0g^2+\cdots$

\textbf{Basis.} $O_S=\tfrac14\operatorname{tr}F^2$ renormalizes multiplicatively: the identity
mixes diagonally, $F\tilde F$ is $P$-odd, and BRST-exact derivatives vanish on physical states.

\textbf{Acid test.} Passes: $0\to0$ (multiplicative); the free $\mathord{:}\varphi^2\mathord{:}$
has $f_1^2=0$ by $\mathbb{Z}_2$ parity ($a\equiv0$); the constant $Z_O^2>0$ because it is a
square, and is $O(1)$ in $V$.

\textbf{Residual.} Conditional on $(\mathrm{CP})$: a single $T=0$ confining phase, a $0^{++}$
edge that is Kato-preserved, and uniform isolation. The theorem
$f_1^2(2.5)\ge Z_O^2r_0(\beta_{\mathrm{eff}})$ is \LC{} on $(\mathrm{CP})$. It transports the
\emph{sign and positivity} of the residue; it transports neither the size nor the isolation.

\subsubsection{Exact variance decomposition (Savvidy-corrected Brascamp--Lieb)}

\textbf{Formula and acid test.} Stated in Chapter~\ref{sec:feshbach}. The decomposition $G''=S-D$
passes the acid test in the strong sense: the free Gaussian has $D=0$, the $k\to0$ mode has
$\operatorname{Var}(\xi)=\infty\Rightarrow S=0$, and $G''=0-0=0$; no false positivity is
manufactured.

\textbf{Residual.} A uniform bound on the fiber defect $D$ at $\beta=2.5$. A measured
$D/S\approx2\%$ was reported \emph{on the gauge-dependent gluon mode}, not the $0^{++}$
(caveat in Appendix~\ref{app:postaudit}). This is H1 in curvature language.

\subsubsection{The Fredholm determinant identity}

\textbf{Formula.}
\begin{equation}
f_1^2=-\lambda_1\operatorname*{Res}_{z=z_1}\mathcal{G}(z)
=-\frac{1}{z_1}\frac{N_O(z_1)}{\mathcal{D}_O'(z_1)},
\label{eq:fredholm}
\end{equation}
with $\mathcal{D}_O(z)=\det(1-zT_O)$,
$\log \mathcal{D}_O=-\sum_k z^k\operatorname{Tr}(T_O^k)/k$, $z_1=e^{am_1}$ the smallest-modulus
zero (so $\lambda_1=e^{-am_1}=1/z_1$ is the leading transfer-matrix eigenvalue), and ``periodic
orbits'' $\operatorname{Tr}T_O^k=Z_{S^1_k\times T^3}$ (closed-time flux tubes).

\textbf{Acid test (strict).} The free field has a pure cut and no pole, so its residue is $0$;
YM has an isolated zero below its own two-glueball cut ($m_1^2<4m_1^2$), so $f_1^2>0$.

\textbf{What it buys.} The pole and the cut have identical lower moments; only the zero set
differs. So the determinant is \textbf{linear-blind} but \textbf{sees the atom}, and it supplies
$m_1^2$ (the zero location), which no moment does.

\textbf{Residual.} Evaluating the zero at $a\to0$ is the marginal BPHZ wall; the continuum
determinant's renormalization \emph{is} the marginal gauge BPHZ, and the zero coalesces into the
K\"all\'en--Lehmann pole (\code{CONTINUUM-FREDHOLM.md}).

\subsubsection{Gauge-invariant IR mass}

\textbf{Formula.}
\begin{equation}
m_{\mathrm{GI}}:=\inf\operatorname{supp}\rho_s
=\lim_{|x|\to\infty}-\partial_{|x|}\log|G_s(x)|=E_1^{(s)}-E_0 ,
\label{eq:mgi}
\end{equation}
on the physical Hilbert space
$\mathcal{H}_{\mathrm{phys}}=L^2(\mathcal{A}/\mathcal{G},\bar\mu)$, with
$s=(\beta/2g)\operatorname{tr}F_{\mu\nu}F^{\mu\nu}\in \mathcal{H}_{\mathrm{phys}}$ and
$\rho_s\ge0$ by reflection positivity.

\textbf{What it removes.} The claim that the gauge potential $A\notin \mathcal{H}_{\mathrm{phys}}$
and that $\rho_A$ has indefinite sign are both artifacts of using $A$; with $F^2$ the spectral
density is non-negative. \textbf{Gauge invariance was never the wall.}

\textbf{Acid test (definitional).} Free Maxwell has $m_{\mathrm{GI}}=0$ (a threshold two-photon
cut from zero); YM has $m_{\mathrm{GI}}=m_{0^{++}}$.

\textbf{Residual.} The edge atom $f_1^2>0$; the definition removes a pseudo-obstruction but does
not manufacture the atom.

\subsubsection{Gauge W-entropy}

\textbf{Formula.}
\begin{equation}
F_W[\mu_t]=\int_{\mathcal{A}}\bigl(S_W(A)+|\nabla_Af_t(A)|_G^2\bigr)e^{-f_t(A)}\,d\mathrm{vol}_G ,
\label{eq:wentropy}
\end{equation}
along the conjugate flow $\partial_tA=-\nabla_AS_W$,
$\partial_tf=-\Delta_Af+|\nabla_Af|^2-S_W$. Gap preservation is monotone iff
$\Gamma_2=\operatorname{Ric}_{\mathcal{A}}+\operatorname{Hess}_{\mathcal{A}}S_W\ge\rho_{\mathrm{Ric}}$.

\textbf{Acid test (mirror).} Free field: $\Gamma_2\ge0$, monotone, but it certifies only the free
lattice gap (non-empty). YM at $\beta=2.5$: $\Gamma_2\approx-2.0<0$ --- correct on the gapless
side, \textbf{blocking} on the gapped side.

\textbf{Verdict.} Perelman's $W$ requires ``geometry $=$ flow'' to be the same object; in
flat-space YM they split. What must be crossed is not a new functional but the missing
non-perturbative flow itself.

\subsubsection{Feshbach / Krein}

The Feshbach--Bogoliubov certificate is Chapter~\ref{sec:feshbach}: an unconditional abstract
two-block lemma whose only debt is TCI. The \textbf{Krein-space} analysis is the precise language
for the indefinite Savvidy Hessian $-D^2+2\operatorname{ad}F_B$: Pontryagin spaces, Langer
definitizable operators, spectral functions. It rigorously proves the UV tower cone-empty
statement (shells $\lambda_j>k_\star$ have $\kappa=0$, $\operatorname{Tr}H_-\lesssim\int|F|^3$)
--- a genuine \textbf{UV half-theorem}. But every Krein invariant is a \textbf{count/sign}
(negative index $r_-$, signature) or a \textbf{UV scale} (the $J$-positive floor
$=\Theta(g|B)$, Landau levels), whereas the gap is the \textbf{min} at the bottom of $k\to0$, an
\textbf{IR} scale. The two are mutually exclusive: choosing $J=\operatorname{sign}H$ makes the
floor physical but gives $\kappa=\dim H_+=\infty$ and is circular (needs a positive spectrum);
choosing an a priori BRST/Gupta--Bleuler metric gives $\kappa=O(1)$ but is literally identical on
the free field (acid-vacuous). A signed negative index is an integer, hence rigid (B3), hence
equal to its free value $0$. The definitizing split loses uniform margin exactly at $k=0$ (the
Savvidy band bottom, where the coarse mode and the most negative mode coincide) --- the Krein
critical point is precisely the point at which the gap must be read. The residual is
$\inf\operatorname{spec}(H|_{\text{physical side}})\ge\kappa_{\mathrm{IR}}>0$, $(L,a)$-uniform,
which is H1.

\subsection{The tool wave: a structural classification}

After the first six tools left their residuals at H1, a dedicated wave built eight
complex-function-theory tools aimed at the identified \emph{forcing blind spot}: a non-integral,
forcing, continuum-robust lower bound on the atom. Seven of the eight fail, and the failure is
structural rather than a tool-choice artifact. The classification is exact
(\code{TOOL-SYNTHESIS.md}):

\begin{enumerate}
\item \textbf{Values / moments / Taylor coefficients / point evaluations of $\Pi$} are finite
  measurable \textbf{integrals} with bounded kernels. The OPE-RHO no-go applies verbatim: they
  are blind to a thin atom and at best give \textbf{upper} bounds ($\Delta_{\max}$, Christoffel,
  Nevanlinna--Pick range).
\item \textbf{Zeros / poles / residues of $\Pi$} (Pad\'e, Wiener--Hopf, argument principle) are
  the atom's only non-integral content; they are reachable only by analytic continuation, i.e.\
  \textbf{A3: read, not force}.
\item \textbf{Positivity} (Carath\'eodory--Fej\'er, Christoffel) \textbf{upper-bounds} an usc
  object; it cannot force a lower bound.
\item \textbf{Geometry / potential / conformal / capacity} are \textbf{support-only}; a point mass
  has zero capacity, so a point atom is invisible to every geometric invariant.
\item \textbf{Convex duality} is \textbf{exact} (zero gap) with the atomless measure feasible, so
  it inherits the blindness: the dual optimum is $0$.
\end{enumerate}

\begin{table}[t]
\centering
\footnotesize
\begin{tabularx}{\textwidth}{@{}>{\raggedright\arraybackslash}p{2.7cm}
  >{\raggedright\arraybackslash}p{1.5cm}>{\raggedright\arraybackslash}X@{}}
\toprule
\textbf{Tool (file)} & \textbf{Verdict} & \textbf{Mechanism}\\
\midrule
Dual / LP (\code{TOOL-DUAL}) & \LN & strong duality gap $=0$; the atomless measure is feasible
  (mass at $M/\chi$), dual optimum $0$ $\Rightarrow$ no dual trick; the blind spot is
  convex-structural\\
Potential (\code{TOOL-POTENTIAL}) & \LN & capacity / equilibrium measure are support-only; a
  point mass has zero capacity; reading the log-coefficient is A3\\
Conformal / extremal (\code{TOOL-CONFORMAL}) & \LN & every conformal invariant is support-only;
  the resolvent level-set radius $\approx f_1^2/R$ is A3\\
Carath\'eodory (\code{TOOL-CARATHEODORY}) & \LN & Taylor coefficients are integrals
  $\int d\nu/(s+z_0)^{k+1}$; positivity is an upper bound only (atom usc)\\
Nevanlinna--Pick (\code{TOOL-NEVANLINNA}) & \LN & point evaluations are integrals (OPE-RHO);
  the LP range $[0,f_{\max}]$ includes $0$; values $\to$ integral class, zeros $\to$ A3\\
Toeplitz / Szeg\H{o} (\code{TOOL-TOEPLITZ}) & \LN & Christoffel is an upper bound
  ($\lambda_n\downarrow f_1^2$); no finite-moment lower bound exists ($\Delta_{\min}=0$)\\
Hardy / Nehari, $L^2$ singular kernel (\code{TOOL-HARDY}) & \textbf{partial} & quadratic
  singular-kernel functional escapes OPE-RHO; detects presence; weight needs location (gap)
  $+\ \varepsilon\to0$ (A3); $L\to\infty$ survival $=$ H1\\
Zero-counting (\code{TOOL-ZEROCOUNT}) & \LN & count is not a faithful level counter; positivity
  admits $s_2=s_1$; cannot exclude isolation $0$; reads $s_1$, not the gap\\
\bottomrule
\end{tabularx}
\caption{The eight complex-function-theory tools. Seven are blind or read-only; one (Hardy)
detects atom presence but not its weight.}
\label{tab:toolwave}
\end{table}

\textbf{Why the classification is a theorem, not a survey.} Each row is a statement about a
\emph{class} of functionals: bounded-kernel integrals are weak-$*$ continuous and hence blind
(Chapter~\ref{ch:weakstar}); zero locations are A3; support functionals cannot see a point; the
dual optimum is zero because the atomless measure is feasible. The wave thus proves that
\emph{no complex-function-theory inequality forces the atom}, which is a structural no-go of
method, not a list of failures.

\subsection{The one genuine escape: the Hardy presence detector}
\label{sec:hardy}

The Hardy functional is
\begin{equation}
N(\varepsilon)=\int_{\varepsilon<|x-m_1^2|<\delta}\bigl|G(x+i0)\bigr|^2\,dx ,
\label{eq:hardy}
\end{equation}
where $G$ is the boundary density (the resolvent/ Stieltjes transform of $\nu$ near the real
axis). It is \textbf{quadratic} in $\nu$ with a \textbf{singular (unbounded) kernel}, so it is
\emph{not} a bounded-kernel integral and OPE-RHO does not apply. The properties are
(\code{TOOL-HARDY.md}, \code{FALSIFY-HARDY-UPGRADE.md}, \code{FALSIFY-HIGHER-SINGULAR.md}):

\begin{itemize}
\item \textbf{Diverges iff atom}, finite iff atomless: as $\varepsilon\to0$,
  $\varepsilon N(\varepsilon)\to2f_1^4$. This is acid-clean: on the free massless field
  ($f_1^2=0$) it is finite; on the gapped-atomless $\mathord{:}\varphi^2\mathord{:}$ it is
  finite; on YM it diverges.
\item \textbf{Continuum-robust}: the recovery coefficient is stable under a background continuum
  (measured $0.49984$, $0.49984$, $0.49987$ for continuum mass $0$, $0.25$, $1.0$), in contrast
  to the raw complex-zero functional, which is destroyed by $0.01\%$ continuum.
\item \textbf{Distinguishes an atom from an equal-mass smooth bump}: the latter is finite at the
  same location, so the divergence is genuinely atom-specific.
\end{itemize}

\textbf{What it does not do.} The quantitative weight requires (i) the pole \textbf{location}
$m_1^2$, whose determination is the gap itself; (ii) the limit $\varepsilon\to0$, which is A3
(a reading across the real axis); and (iii) the $L\to\infty$ survival of the isolated pole, which
is isolation I2 $=$ H1. The singular-kernel route cannot decouple ``location $+$ $\varepsilon\to0$''
from the atom: they are the same datum. Pushing to higher singularity
($\int|G^{(k)}|^2$, $\int|G|^p$) only changes the divergence power $\varepsilon^{-(p-1)}$, not
the need for the location and the limit.

\textbf{Why this is still valuable.} It is the first tool in the campaign that genuinely escapes
the bounded-kernel integral no-go \emph{for presence}. It converts ``does the atom exist?'' from
a blind question to a computable one, at the price of not computing the weight. In the
two-coordinate language of Part~III, it is a weak-$*$-discontinuous functional that is
nevertheless computable from a finite window --- the one exception that sharpens the rule.

\subsection{Scorecard}

\begin{table}[t]
\centering
\footnotesize
\begin{tabularx}{\textwidth}{@{}>{\raggedright\arraybackslash}p{3.0cm}
  >{\raggedright\arraybackslash}p{2.9cm}c
  >{\raggedright\arraybackslash}p{2.6cm}
  >{\raggedright\arraybackslash}X@{}}
\toprule
\textbf{Tool} & \textbf{Formula / object} & \textbf{Acid test} & \textbf{Touches the atom?} &
\textbf{Exact residual}\\
\midrule
Multiplicative transport & $f_1^2(\mu)=Z_O^2f_1^2(\mu_0)$ & pass ($0\to0$, square) &
  sign/size transport & $(\mathrm{CP})$\\
Variance decomposition & $G''=S-D$ & pass (free $0-0$) & curvature sign &
  uniform $D$ at 2.5 $=$ H1\\
Fredholm determinant & $f_1^2=-\lambda_1\operatorname{Res}G$ & strict pass &
  \textbf{yes} (zero set) & $a\to0$ zero\\
Gauge-invariant IR mass & $m_{\mathrm{GI}}=\inf\operatorname{supp}\rho_s$ &
  pass (Maxwell 0) & removes artifact & edge atom\\
Gauge W-entropy & $F_W$, $\Gamma_2\ge\rho_{\mathrm{Ric}}$ & mirror pass & no (gap side blocks) &
  missing flow\\
Feshbach/Krein & two-block lemma, $\Pi_\kappa$ & pass (free returns 0) &
  reduction to TCI & TCI $=$ H1\\
Hardy presence & $N(\varepsilon)=\int|G|^2$ & strict pass & presence only &
  location (gap) $+$ $\varepsilon\to0$ (A3)\\
Seven CF tools & values / zeros / support / dual & various & no & integral class or A3\\
\bottomrule
\end{tabularx}
\caption{Scorecard for the six first-wave tools and the tool wave.}
\label{tab:tools}
\end{table}

Four of the tools --- transport, variance decomposition, determinant, Hardy --- genuinely
``touch'' the atom; all four leave the wall on the far side of ``continuum plus transport to
2.5,'' which is the analyticity/BPHZ wall, or of the location-plus-limit problem, which is A3.
The pattern is the same in every column, and Chapter~\ref{ch:weakstar} explains why it must be.
